\section{The Gelman-Rubin $\hat{R}$ Statistic}
\label{sec:rhat}

\subsection{Definition}

Let $f : \mathcal{P} \to \mathbb{R}$ be a scalar function of
redistricting plans (e.g., Polsby-Popper compactness or Democratic
seat count).
Run $m$ independent \recom{} chains of length $n$ each, and denote
the $j$-th draw from chain $i$ as $\psi_{ij} = f(\pi_{ij})$.

Define the between-chain variance:
\[
  B = \frac{n}{m-1} \sum_{i=1}^m \left(\bar{\psi}_{i\cdot} - \bar{\psi}_{\cdot\cdot}\right)^2,
\]
and the within-chain variance:
\[
  W = \frac{1}{m} \sum_{i=1}^m s_i^2, \quad
  s_i^2 = \frac{1}{n-1}\sum_{j=1}^n \left(\psi_{ij} - \bar{\psi}_{i\cdot}\right)^2.
\]
The marginal posterior variance estimator is:
\[
  \widehat{\mathrm{var}}^+(\psi \mid y) = \frac{n-1}{n} W + \frac{1}{n} B.
\]
The potential scale reduction factor is:
\[
  \rhat = \sqrt{\frac{\widehat{\mathrm{var}}^+(\psi \mid y)}{W}}.
\]

If the chains have not mixed, $B \gg W$ and $\rhat \gg 1$.
If the chains have converged to the same stationary distribution,
$B / n \approx W$ and $\rhat \approx 1$.

\subsection{Threshold and Interpretation}

Following \citet{gelman1992} and the updated analysis of \citet{vehtari2021},
we adopt the threshold $\rhat < 1.1$ for approximate convergence.
This threshold is conservative: it permits $\widehat{\mathrm{var}}^+$ to
exceed $W$ by up to 10\%, implying the chains have not fully decorrelated
but have explored broadly similar regions.

For redistricting applications, we compute $\rhat$ for two outcome
functions:
\begin{enumerate}
\item \textbf{State-average Polsby-Popper ($\PP$)}: the mean of the
      district-level Polsby-Popper scores within a plan.
\item \textbf{Democratic seat count ($D$-seats)}: the number of districts
      with Democratic two-party vote share exceeding 50\%.
\end{enumerate}

$\PP$ converges faster than $D$-seats because compactness is a continuous
function of district geometry and varies smoothly with small plan changes.
$D$-seats is integer-valued and can jump discontinuously when a swing
district crosses the 50\% threshold.

\subsection{Application to Ten ReCom Chains}

For each of the six study states, we run $m = 10$ independent \recom{}
chains starting from randomised initial plans (generated by random
spanning-tree bisection with no compactness objective).
Each chain runs for $n = 10{,}000$ steps.

We compute $\rhat$ at checkpoints $n \in \{500, 1{,}000, 2{,}000, 5{,}000, 10{,}000\}$
to track convergence progress.
Section~\ref{sec:application} presents the full results.
Typical findings: $\rhat(\PP) < 1.1$ is achieved at $n \approx 2{,}000$
steps for small states (NC, WI) and $n \approx 10{,}000$--$20{,}000$
steps for large states (TX, CA).

\subsection{Rank-Normalized $\hat{R}$}

\citet{vehtari2021} recommend rank-normalizing the draws before computing
$\rhat$ to improve robustness against heavy-tailed distributions.
For redistricting outcomes, Polsby-Popper is bounded in $[0, 1]$ and
approximately Beta-distributed, making the original $\rhat$ reliable.
Democratic seat counts are bounded integers; we apply rank normalization
for this variable.
Numerical results in Section~\ref{sec:application} report both the
original and rank-normalized $\rhat$ for comparison.
